\documentclass[11pt,a4paper]{article}
\usepackage[margin=25mm]{geometry}
\usepackage{amsmath,amssymb,amsthm,hyperref}
\setlength{\emergencystretch}{2em}
\newtheorem{lemma}{Lemma}
\newtheorem{proposition}[lemma]{Proposition}
\newtheorem{corollary}[lemma]{Corollary}
\newcommand{\gen}[1]{\langle#1\rangle}
\title{Completing the numerical intransitive-side support bound}
\author{Human Compression project: bounded L044 proof component}
\date{}
\begin{document}\maketitle
Throughout, $G=\gen{\rho_0,\ldots,\rho_{r-1}}\leq A_n$ is
transitive, $n\geq12$, its marking is a connected string C-marking,
and every maximal omission is intransitive. The unique eligible
interior $i$-edge $\{a,b\}$ joins the two $G_i$-orbits $A,B$.
Assume explicitly that $|B|\geq2$ and that $L=G_{<i}$ fixes $B$
pointwise and is intransitive on $A$. Put $R=G_{>i}$ and
$\tau=\rho_{i+1}$. Write $m$ for the number of $L$-orbits in $A$
and $k$ for their common size.

We use the already proved faithful commuting-orbit facts: $L$ is
faithful on each such orbit, $k\geq i+1$, and $\tau$ exchanges
$\beta=a^L$ with a distinct orbit $\beta'$. Every generator of
index at least $i+2$ fixes $\beta$ pointwise, and those of index
at least $i+3$ fix $\beta'$ pointwise. The only generator of $G_i$
which can move $b$ is $\tau$; since $B$ is nontrivial, it does.
These facts use the unique-crossing lemma, commuting factors and
Whiston's independent-length bound, not IP of a projected action.

\begin{lemma}[Distance to a moved block]\label{blocks}
If the last generator has nontrivial image on the $L$-orbits in
$A$, and $d=r-1-i\geq2$, then $m\geq d+1$.
\end{lemma}
\begin{proof}
In the block graph, take a shortest path from $\beta$, using
indices $i+1,\ldots,r-2$, to a block moved by $\rho_{r-1}$.
Such a block is reachable: otherwise the set of reachable blocks
would be invariant under every generator of $R$, hence would be
all blocks by transitivity, contradicting the assumed nontrivial
last block image.
Every intermediate index occurs in this path word. If index $j$
were missing, the factors below and above $j$ could be reordered,
since they commute. The upper factors fix $\beta$; the lower
factors commute with $\rho_{r-1}$, which fixes $\beta$ as well.
The endpoint block would then be fixed by $\rho_{r-1}$, a
contradiction. The path therefore has at least $d-1$ edges and
$d$ distinct blocks. All earlier blocks are fixed by the last
generator, by minimality. The image of the last block under that
generator is a further block, proving the bound.
\end{proof}

\begin{proposition}[A final generator supported on the first orbit]
\label{first}
If $\rho_{r-1}$ fixes $B$ pointwise, then
$r\leq\lfloor(n-1)/2\rfloor$.
\end{proposition}
\begin{proof}
The last generator is nonidentity and supported on $A$.
Its block image cannot belong to the group generated by the
other $R$ block images. If it did, those other images would
remain transitive on the blocks. Together with $L$ they would
make the global last omission transitive on $A$. Adjoining a
generator supported on that one omission orbit could not join
any omission orbits. Since $G$ is transitive, the omission would
already be transitive, a contradiction. In particular the last
block image is nontrivial. Also $d=r-1-i\geq2$, because $\tau$
moves $B$. Lemma~\ref{blocks} gives
\[
 n=km+|B|\geq(i+1)(d+1)+2
       =2r+(i-1)(d-1).
\]
For $i\geq2$ this is at least $2r+1$, as required.

Suppose $i=1$ and the rank bound fails. All the inequalities
must then be equalities:
\[
 k=2,\qquad |B|=2,\qquad m=d+1,\qquad n=2r=2m+2.
\]
The generator $\rho_0$ exchanges every one of the $m$ pairs in
$A$ and fixes $B$. It is even, so $m$ is even; $n\geq12$ gives
$m\geq6$.

Every distinguished $R$ block image is essential. The image of
$\tau$ is essential because all the others fix $\beta$.
Each other generator fixes $b$, hence all of the two-point set
$B$. If its block image were dispensable, the same global
omission argument as above would apply to that generator.
Thus the $d=m-1$ block images are an independent involutory
string in $S_m$. Whiston's equality theorem makes the image
$S_m$. The paid maximal symmetric-string classification gives
the ordinary adjacent-transposition path, except possibly the
outer degree-six Coxeter action. The latter has transitive
endpoint omissions, whereas omitting $\tau$ here fixes
$\beta$, so it is excluded. This uses the accepted complete
degree-six certificate and the uniform theorem for $m\geq7$;
it assumes no projected IP.

Number the pairs $0,\ldots,m-1$ so the successive $R$ block
images swap $0,1$, then $1,2$, and so on. The anchor pair is
pair $0$: it is the fixed point of the block subgroup omitting
$\tau$. Put $H=\gen{\rho_{i+3},\ldots,\rho_{r-1}}$.
This group fixes the first two pairs pointwise and induces
$S_{m-2}$ on the remaining $m-2$ pairs. On the first four
points, $\tau$ exchanges two pairs, so is a product of two
transpositions. On the remaining points it centralizes $H$.
The accepted paired-row centralizer lemma applies because
$m-2$ is even and at least four: any such centralizing
involution is even. Hence $\tau$ is even on $A$ but is a
single transposition on $B$. This contradicts $G\leq A_n$.
\end{proof}

\begin{proposition}[The last generator shared by both orbits]
\label{last}
If $\rho_{r-1}$ acts nontrivially on both $A$ and $B$, then
$r\leq\lfloor(n-1)/2\rfloor$.
\end{proposition}
\begin{proof}
Put $d=r-1-i$. For $d\geq2$, the accepted active-orbit distance
lemma gives $|A|\geq kd$. The fixed-prefix/no-gap lemma, applied
in $B$ and including the final edge, gives a simple path with
at least $d+1$ vertices there. Consequently
\[
 n\geq d(k+1)+1\geq d(i+2)+1
          =2r+i(d-2)-1.                                  \tag{1}
\]
If $i\geq2$ and $d\geq3$, this is at least $2r+1$.
If $i=1$ and $d\geq4$, it is again at least $2r+1$.
The remaining case $i=1,d=3$ has $r=5$, and the bound follows
directly from $n\geq12$.

For $d=2$, the first orbit on which the last generator can
act is $\beta'$. If it moves $\beta'$ to another orbit,
the active-orbit lemma gives $|A|\geq3k$, while $|B|\geq3$.
Thus $n\geq3k+3\geq3i+6>2i+6=2r$.
If it preserves $\beta'$ nontrivially, the complete
internal-block lemma applies and gives the asserted terminal
rank bound. It cannot fix $\beta'$ pointwise: then
$\beta\cup\beta'$ would be invariant under all of $R$,
making it all of $A$ and leaving the last generator trivial
there. This accounts for every $d=2$ possibility.

It remains to consider $d=1$. Now $R=\gen\tau$, so
$A=\beta\cup\beta'$ and $B=\{b,b'\}$.
Thus $n=2k+2$ and $r=i+2\leq k+1$. Failure of the rank bound
would force $i=k-1$ and $r=k+1$. The involution $\tau$
exchanges the two $k$-point orbits and also exchanges $b,b'$;
its evenness forces $k$ odd, and $n\geq12$ gives $k\geq5$.

The faithful degree-$k$ action of $L$ has an independent string
of length $k-1$. Whiston and the paid maximal symmetric-string
classification identify it as the ordinary path marking of
$S_k$; the degree-five certificate covers $k=5$, and the
uniform result covers odd $k\geq7$.
The subgroup $H=G_{<i-1}$ fixes $a$ and $a'=a^\tau$ and
acts diagonally as the natural $S_{k-1}$ on the remaining
$k-1$ pairs, one point from each of $\beta,\beta'$.
These are its only two fixed points in $A$.

Write the crossing involution as $s=\rho_i=(a\ b)u$.
The permutation $u$ is supported on $A\setminus\{a\}$:
the cross-edge is unique, and the other point $b'$ of $B$
cannot be moved. The group $H$ fixes $a,b$ and commutes with
$s$, so $u$ centralizes $H$. It fixes $a$, hence fixes the
other $H$-fixed point $a'$ as well. On the remaining
$2(k-1)$ points the paired-row centralizer lemma applies,
since $k-1$ is even and at least four. Thus $u$ is even,
making $s$ odd, a contradiction. This completes the last case.
\end{proof}

\begin{corollary}[Complete numerical support alternative]
Let $h$ be the largest index acting nontrivially on both $A$
and $B$. Either $r\leq\lfloor(n-1)/2\rfloor$, or $h<r-1$
and, writing $X=\operatorname{supp}(G_{>h})$, one has
\[
 X\subseteq B\setminus\{b\},\qquad
 n-|X|\geq2h+1.
\]
The latter inequality is strict if $h\geq i+2$. At equality,
$h=i+1$, $k=i+1$, $A=\beta\cup\beta'$, the fixed set is
exactly $A\cup\{b\}$, and the faithful action of $L$ on
either of its two orbits is $S_k$.
\end{corollary}
\begin{proof}
The index $h$ exists because $\tau$ moves both $A$ and $B$.
If the rank bound fails, Proposition~\ref{last} excludes
$h=r-1$. Every later generator is nonidentity and acts on
exactly one of the two orbits. Connectedness of the string
diagram forces all these generators to act on the same
orbit: a change of sides would make an adjacent pair commute
and disconnect the diagram. Proposition~\ref{first} excludes
$A$ as that side. Thus the tail fixes $A$ pointwise, and the
accepted orbit-distance support proposition gives every
remaining assertion, including the strictness and equality
parameters.
\end{proof}

\paragraph{Source, paid inputs and precise scope.}
This completes the numerical support alternative in the
intransitive-left-side branch of Cameron--Fernandes--Leemans--Mixer,
\emph{Highest rank of a polytope for $A_n$},
arXiv:1605.09173v1, Proposition~5.23, physical pages~32--33.
The source's commuting-block framework is retained; the orbit
distance and tight-case parity arguments above are the present
replacement proof. The exact fixed-set conclusion does not
assert the source's complete coloured equality graph.
All unique-crossing, active-orbit, internal-block and paired-row
lemmas remain paid dependencies, as do the full L047 maximal
symmetric-string theorem, its complete small certificate and
the retained Whiston bounds. The internal-block input is its
explicit nontrivial-$B$ v2; the earlier frozen v1 is preserved
historically. No new finite computation or formal declaration
is needed. Singleton $B$, the residual small attachment and
full-IP equality rigidity, later source tables, and the whole
L044 contract remain separate. There is no new original,
question or priority claim.
\end{document}
